\chapter{Validation and errors}

Validation is the first line of defense.
If you accept bad input, every downstream system pays the cost.

\section{Validation strategy}
Validate at the edge and again at the core for critical invariants.
Use strict schemas for external requests.

\section{Error shape}
Return one error format to clients and log a richer internal error.

\begin{lstlisting}[language=JSON]
{
  "error": {
    "code": "invalid_request",
    "message": "title is required",
    "field": "title"
  }
}
\end{lstlisting}

\section{Field errors}
Expose field errors for forms and batch operations.
Keep error codes stable across versions.

\section{Mapping errors to status codes}
Validation errors should be \code{400}.
Permission errors should be \code{403}.
Missing resources should be \code{404}.

\section{Where validation lives}
Validation happens at multiple layers:
\begin{itemize}[nosep]
  \item edge validation for request shape
  \item domain validation for business rules
  \item database constraints for invariants
\end{itemize}

If any layer is missing, invalid data will leak through.

\section{Schema validation}
Use explicit schemas for external inputs.
Schema validation should reject unknown fields when you want strict contracts.
If you allow unknown fields, document that behavior clearly.

\section{Partial update validation}
For \code{PATCH}, validate only the fields provided.
Do not allow partial updates to break invariants.
If a field depends on another field, validate that dependency explicitly.

\section{Error taxonomy}
Keep a small, stable set of error codes.

\begin{tabularx}{\textwidth}{@{}lX@{}}
\toprule
Code & Meaning \\
\midrule
invalid\_request & request payload is malformed or missing fields \\
validation\_error & input fails business rule validation \\
not\_found & resource does not exist or is not visible \\
conflict & request violates a uniqueness or state constraint \\
rate\_limited & request exceeded a rate limit \\
\bottomrule
\end{tabularx}

\section{Batch validation}
Batch operations should return per-item errors.
Do not fail an entire batch when only one item is invalid.

\begin{lstlisting}[language=JSON]
{
  "results": [
    {"id": "task_1", "status": "created"},
    {"id": null, "status": "error", "error": "validation_error"}
  ]
}
\end{lstlisting}

\section{Localization and messages}
Error messages are for humans, error codes are for machines.
If you localize messages, keep the codes stable across languages.

\section{Logging and observability}
Log validation failures with a consistent structure.
This makes it easy to detect which fields cause the most errors.
Do not log sensitive values such as passwords or tokens.

\section{Retry guidance}
Validation errors are not retryable.
Conflict errors may be retryable with updated input.
Rate limit errors should include \code{Retry-After}.

\section{Example validation flow}
For \code{POST /projects/:id/tasks}:
\begin{itemize}[nosep]
  \item validate JSON shape and required fields
  \item confirm project exists and membership is valid
  \item enforce status and date rules
  \item return a structured error on failure
\end{itemize}

\section{Constraint violations}
Database constraints are a final safety net.
Map constraint errors to stable error codes.
Do not leak internal table or column names to clients.

\section{Consistency with logging}
Client errors should be logged at a lower level.
Server errors should be logged at error level with full context.
Avoid logging full request bodies for privacy reasons.

\section{Validation for imports}
Bulk imports need different validation.
Provide an error report with row numbers and field errors.
This helps users fix data without guessing.

\section{Normalization}
Normalize inputs before validation when possible.
Trim whitespace, normalize case, and parse dates consistently.
This reduces surprising validation failures.

\section{Reserved values}
Some values are reserved or system-defined.
Validate them explicitly:
\begin{itemize}[nosep]
  \item prevent reserved names
  \item block forbidden status transitions
  \item enforce unique identifiers
\end{itemize}

\section{Error metadata}
Error responses should include a request id.
If a support ticket comes in, that id should locate the error in logs.

\section{Example error response with metadata}
\begin{lstlisting}[language=JSON]
{
  "error": {
    "code": "validation_error",
    "message": "status is invalid",
    "field": "status",
    "request_id": "req_42"
  }
}
\end{lstlisting}

\section{Validation in async jobs}
Background jobs should validate inputs too.
If a job receives invalid data, fail the job and surface the error.
Silent failures create data drift.

\section{Client-side validation}
Client-side validation improves UX but is not a security boundary.
Server-side validation is the source of truth.

\section{Idempotency and validation}
If a request is retried, validation should be consistent.
Do not allow a retried request to pass validation when the first one failed.

\section{Error contracts}
Treat error shapes as part of the API contract.
Changing them can break clients just as much as changing fields.

\section{Validation as a product surface}
Users experience validation failures before they see your core features.
Clear, precise error messages reduce support load and increase adoption.
Treat validation text as part of your product design.

\section{Field path conventions}
Standardize field paths so clients can map them to UI elements.
Use dot-notation or JSON pointer and keep it consistent across endpoints.
Do not invent new field path formats per API.

\section{Partial update rules}
PATCH endpoints need explicit rules about missing fields.
Define which fields can be cleared, which are immutable, and which require validation only when present.
Ambiguity here is a common source of regressions.

\section{Example partial update policy}
\begin{lstlisting}[language=JSON]
{
  "update_policy": {
    "title": "optional",
    "status": "optional",
    "owner_id": "immutable",
    "due_date": "nullable"
  }
}
\end{lstlisting}

\section{Batch validation details}
Batch APIs should return per-item errors with stable indexes.
Do not stop at the first failure unless the operation is transactional.
Return a summary error and per-item details.

\section{Error category semantics}
Clients need to know if an error is fixable.
Separate errors into validation, authorization, conflict, and system categories.
Expose a stable \code{type} value for each category.

\section{Error category table}
\begin{tabularx}{\textwidth}{@{}lX@{}}
\toprule
Category & Client action \\
\midrule
Validation & Fix input and retry \\
Authorization & Request access or credentials \\
Conflict & Refresh state and retry if still relevant \\
System & Retry with backoff, or contact support \\
\bottomrule
\end{tabularx}

\section{Error code ownership}
Assign ownership of error codes to modules.
If everyone can mint codes, you end up with duplicates.
Ownership also helps with documentation and deprecation.

\section{Localization strategy}
Decide whether servers return localized messages or message keys.
Message keys keep the API stable and let clients localize independently.
If you localize server-side, document supported locales and fallbacks.

\section{Normalization and preprocessing}
Normalize inputs before validation.
Trim whitespace, normalize Unicode, and coerce types where appropriate.
Validation without normalization leads to inconsistent results.

\section{Reserved values revisited}
Reserve values like \code{admin} or \code{root} explicitly.
Do not rely on implicit checks scattered across the codebase.
Centralize reserved value checks in one validation module.

\section{Constraint evolution}
Constraints change over time.
Introduce new constraints in warning mode before enforcing them.
This avoids breaking existing clients without notice.

\section{Async validation}
Some validations require external calls.
Separate async validation from synchronous checks.
Return a distinct error type so clients can handle the difference.

\section{Logging and privacy}
Validation errors can include sensitive data.
Log only sanitized versions of fields.
Avoid logging full payloads in production.

\section{Retries and idempotency}
Some errors are safe to retry, others are not.
Include a \code{retryable} flag when relevant.
If retries are allowed, ensure the operation is idempotent.

\section{Example error envelope}
\begin{lstlisting}[language=JSON]
{
  "error": {
    "type": "validation",
    "code": "task.title.too_long",
    "message": "Title must be 80 characters or fewer",
    "field": "title",
    "retryable": false
  }
}
\end{lstlisting}

\section{Auditability}
Validation rules are part of system policy.
Document them and version them.
This helps audits and reduces ambiguity in regulated environments.

\section{Testing validation}
Tests should cover boundary values, not just happy paths.
Include property-based tests for input ranges where possible.
Regression tests for validation are cheap and valuable.

\section{Validation architecture}
Keep validation rules close to the domain model.
Transport layers should handle parsing, but core validation should live with business logic.
This prevents duplicate logic across HTTP, CLI, and background jobs.

\section{Strict versus lenient modes}
Some endpoints need strict validation, others tolerate extra fields.
Define which mode each endpoint uses.
Switching modes silently is a breaking change.

\section{Backward compatibility}
Changing validation rules can be a breaking change.
If you tighten rules, consider a deprecation window.
Provide guidance to clients before enforcement.

\section{Error documentation}
Document error codes and fields.
If clients need to guess, they will implement brittle logic.
Error docs are part of your API surface.

\section{CSV import example}
CSV imports need row-level feedback.
Return a summary and a per-row error list with line numbers.
This is often more useful than a single failure message.

\section{Example import error report}
\begin{lstlisting}[language=JSON]
{
  "summary": {"processed": 200, "errors": 12},
  "errors": [
    {"row": 14, "field": "email", "code": "invalid_format"},
    {"row": 87, "field": "status", "code": "reserved_value"}
  ]
}
\end{lstlisting}

\section{Constraint mapping}
Map database errors to user-facing codes.
Do not expose SQL errors or internal constraint names.
Keep a mapping table per database.

\section{Example constraint mapping}
\begin{tabularx}{\textwidth}{@{}lX@{}}
\toprule
Constraint & Public code \\
\midrule
users\_email\_unique & conflict.email\_exists \\
tasks\_due\_date\_check & validation.due\_date\_range \\
\bottomrule
\end{tabularx}

\section{Security validation}
Validate inputs that affect queries or templates.
Reject unexpected characters where appropriate.
Sanitize output even when validation passes.

\section{Field metadata}
Include metadata for client UX, such as max length or allowed values.
Metadata reduces guesswork in clients.
Keep metadata versioned to avoid mismatches.

\section{Error aggregation}
Aggregate multiple errors instead of returning one at a time.
This improves client efficiency and user experience.
Limit the error count to avoid huge payloads.

\section{Validation pipeline}
A simple pipeline keeps logic predictable:
\begin{itemize}[nosep]
  \item normalize input
  \item run schema validation
  \item run business validation
  \item enforce persistence constraints
\end{itemize}

\section{Rate limit errors}
Rate limits are a validation boundary.
Return \code{429} with a clear retry time.
Expose a stable error code for rate limit failures.

\section{Client UX guidance}
If possible, include human-friendly messages.
Messages should be short and specific.
Avoid blaming language such as \emph{invalid} when you can explain the fix.

\section{State transition validation}
Many fields have allowed state transitions.
Validate transitions explicitly rather than relying on ad-hoc checks.
State validation prevents invalid workflows from entering the system.

\section{Warning responses}
Sometimes you want to accept input but warn about future behavior.
Provide warnings separately from errors.
Warnings prepare clients for upcoming constraint changes.

\section{Example warning response}
\begin{lstlisting}[language=JSON]
{
  "warnings": [
    {"code": "deprecated_field", "field": "legacy_id"}
  ]
}
\end{lstlisting}

\section{Help links}
Include a stable help link for common error codes.
Help links reduce support burden.
Keep links versioned when the documentation changes.

\section{Limits and caps}
Validation includes limits like maximum list size or payload size.
Return clear errors when limits are hit.
Do not silently truncate user input.

\section{Contract tests}
Add contract tests for validation responses.
These tests verify that error shapes and codes remain stable.
Contract tests prevent accidental breaking changes.

\section{Numeric precision}
Validate numeric precision explicitly.
Money values should use integer cents or decimal types.
Floating-point errors create subtle data issues.

\section{Time zones}
Date validation must account for time zones.
Store timestamps in UTC and validate input offsets.
Reject ambiguous local times during daylight transitions.

\section{Schema versioning}
When schemas evolve, version them explicitly.
Support multiple schema versions during migrations.
Document the version negotiation strategy for clients.

\section{Error quotas}
Limit the number of errors returned for large payloads.
Provide a count of suppressed errors to signal incomplete feedback.
This keeps error responses fast and predictable.

\section{Example error quota}
\begin{lstlisting}[language=JSON]
{
  "error_count": 500,
  "returned_errors": 100,
  "suppressed_errors": 400
}
\end{lstlisting}

\begin{checklistbox}{Checklist: validation}
\begin{itemize}[nosep]
  \item Validate at the edge.
  \item Keep error shapes consistent.
  \item Use stable error codes for clients.
\end{itemize}
\end{checklistbox}
